\section{Introduction}\label{sec:introduction}

Coding agents increasingly work as orchestrators. A top-level session decomposes a request, hands independent pieces to sub-agents, and merges what comes back. Each hand-off forces a choice: which model should run the sub-agent, and at what reasoning effort? These two settings determine what a delegated slice costs, how long the orchestrator waits, and whether the returned answer can be trusted. A cheap setting that fails wastes a round trip. An expensive setting that was not needed wastes money and time on every slice.

In practice, we expect the choice is often made by habit or by a rule of thumb; a common one is to use the largest model for anything difficult and the smallest for lookups. Effort levels add a second axis with little guidance. Higher effort means more thinking tokens, and one might expect it to help on harder work, but it also raises cost and latency, and it is not obvious that it helps on the tasks sub-agents are actually given. We could not find measurements that let an orchestrator pick a model and effort for a given slice of work.

We study this question directly. We run Claude Haiku~5.5 and Claude Sonnet~5.5 as sub-agents at up to five effort levels, \texttt{low}, \texttt{medium}, \texttt{high}, \texttt{xhigh}, and \texttt{max}. We call one model at one effort level a \emph{config}. Each sub-agent runs as a separate Claude Code session in its own terminal pane, so the setting resembles real orchestration rather than a bare API call. We use three experiments. The first asks agents to reason exactly, without tools, over families of tasks whose difficulty we can scale. The second and third give agents coding tasks in a private Python repository: locating code, fixing planted bugs, replaying real feature commits, and reviewing diffs with planted bugs. Fix and feature tasks are graded by running tests after the agent exits, not by trusting the agent's own report. For every run we compute the API-equivalent cost from the session transcript, so we can compare configs by dollars per successful task as well as by pass rate.

The picture that emerges is mixed. On coding, almost every config passed almost every task, so effort and model choice separated configs mainly through cost and time. On long exact reasoning without tools, they separated configs sharply. A few code-review misses repeated across every effort level of one model, which suggests that a second model may add more coverage than more effort. We say plainly where the evidence is thin: the coding tasks were mostly passed by every config, so the study does not locate the point where coding reliability breaks down.

\paragraph{Contributions.}
\begin{enumerate}[leftmargin=*]
  \item A harness that runs Claude Code sub-agents in terminal panes with controlled permissions, fresh worktrees, and post-hoc grading from the working tree, together with the validation steps we used to check tasks, briefs, and answer keys.
  \item Measurements of two models across effort levels on 229 tool-free reasoning runs and 270 coding runs, reported as pass rates with Wilson intervals and as API-equivalent dollars and agent time per successful task.
  \item A routing table derived from these measurements, with each rule marked as measured, observed in a handful of cases, untested, or a judgment, so readers can see how much weight it can bear.
  \item An explicit statistical framework for a larger follow-up, including an item response model and an error-correlation analysis, and a list of the harness and benchmark-design problems we met along the way.
\end{enumerate}

\paragraph{Roadmap.}
Section~\ref{sec:related-work} reviews related work. Section~\ref{sec:method} describes the harness, the configs, the tasks and their validation, the grading rules, and the cost accounting. Section~\ref{sec:experiments} reports the three experiments and the cost results. Section~\ref{sec:discussion} turns them into a routing table, sets out the modeling framework, discusses a second model as an alternative to more effort, and states lessons, ethics, and limitations. Section~\ref{sec:conclusion} concludes with future work.
